% Options for packages loaded elsewhere
\PassOptionsToPackage{unicode}{hyperref}
\PassOptionsToPackage{hyphens}{url}
%
\documentclass[
]{article}
\usepackage{amsmath,amssymb}
\usepackage{iftex}
\ifPDFTeX
  \usepackage[T1]{fontenc}
  \usepackage[utf8]{inputenc}
  \usepackage{textcomp} % provide euro and other symbols
\else % if luatex or xetex
  \usepackage{unicode-math} % this also loads fontspec
  \defaultfontfeatures{Scale=MatchLowercase}
  \defaultfontfeatures[\rmfamily]{Ligatures=TeX,Scale=1}
\fi
\usepackage{lmodern}
\ifPDFTeX\else
  % xetex/luatex font selection
\fi
% Use upquote if available, for straight quotes in verbatim environments
\IfFileExists{upquote.sty}{\usepackage{upquote}}{}
\IfFileExists{microtype.sty}{% use microtype if available
  \usepackage[]{microtype}
  \UseMicrotypeSet[protrusion]{basicmath} % disable protrusion for tt fonts
}{}
\makeatletter
\@ifundefined{KOMAClassName}{% if non-KOMA class
  \IfFileExists{parskip.sty}{%
    \usepackage{parskip}
  }{% else
    \setlength{\parindent}{0pt}
    \setlength{\parskip}{6pt plus 2pt minus 1pt}}
}{% if KOMA class
  \KOMAoptions{parskip=half}}
\makeatother
\usepackage{xcolor}
\usepackage[left=2cm,right=1cm,top=1cm,bottom=1.5cm]{geometry}
\usepackage{graphicx}
\makeatletter
\def\maxwidth{\ifdim\Gin@nat@width>\linewidth\linewidth\else\Gin@nat@width\fi}
\def\maxheight{\ifdim\Gin@nat@height>\textheight\textheight\else\Gin@nat@height\fi}
\makeatother
% Scale images if necessary, so that they will not overflow the page
% margins by default, and it is still possible to overwrite the defaults
% using explicit options in \includegraphics[width, height, ...]{}
\setkeys{Gin}{width=\maxwidth,height=\maxheight,keepaspectratio}
% Set default figure placement to htbp
\makeatletter
\def\fps@figure{htbp}
\makeatother
\setlength{\emergencystretch}{3em} % prevent overfull lines
\providecommand{\tightlist}{%
  \setlength{\itemsep}{0pt}\setlength{\parskip}{0pt}}
\setcounter{secnumdepth}{-\maxdimen} % remove section numbering
\usepackage{blindtext}
\usepackage{amsmath}
\usepackage{mathtools}
\usepackage{multicol}
\ifLuaTeX
  \usepackage{selnolig}  % disable illegal ligatures
\fi
\IfFileExists{bookmark.sty}{\usepackage{bookmark}}{\usepackage{hyperref}}
\IfFileExists{xurl.sty}{\usepackage{xurl}}{} % add URL line breaks if available
\urlstyle{same}
\hypersetup{
  pdftitle={Tarea Vectores, rectas y planos},
  pdfauthor={álgebra lineal},
  hidelinks,
  pdfcreator={LaTeX via pandoc}}

\title{Tarea Vectores, rectas y planos}
\author{álgebra lineal}
\date{2024-09-14}

\begin{document}
\maketitle

\hypertarget{vectores}{%
\section{Vectores}\label{vectores}}

\begin{enumerate}
\def\labelenumi{\arabic{enumi}.}
\tightlist
\item
  Una partícula tiene una posición en el tiempo de acuerdo a un
  movimiento rectilíneo uniforme dado por
  \(\vec{\mathbf{r}}(t) = 6t\hat{\mathbf{\iota}} + (8t+2)\hat{\mathbf{\jmath}}+(2t+1)\hat{\mathbf{k}}\).
  Suponga que a partir del tiempo \(t=0\) se observa su desplazamiento
\end{enumerate}

\begin{enumerate}
\def\labelenumi{\alph{enumi})}
\item
  ¿Cúal es la posición de la partícula después de 4 segundos?
\item
  Calcular a que distancia del origen se encuentra la partícula al
  segundo 5.
\item
  Si la rápidez de la partícula se calcula como
  \(|\vec{v}|=\sqrt{v_x^2+v_y^{2}+v_z^{2}}\). Calcular la rapidez que
  tiene la partícula en el instante \(t=4\).
\end{enumerate}

\begin{enumerate}
\def\labelenumi{\arabic{enumi}.}
\item
  Investigar el concepto de \emph{cosenos} \emph{directores}. Calcular
  los cosenos directores de los siguientes vectores
  \(\vec{\mathbf{a}}=(12,-15,-16)\), \(\vec{\mathbf{a}} = (2,-3,-1)\) y
  \(\vec{\mathbf{a}} = (3, 4, 5)\)
\item
  Sean \(\vec{\mathbf{a}}=(3,-1,2,4)\), \(\vec{\mathbf{b}}=(1,2,0,1)\),
  \(\vec{\mathbf{c}}=(3,2,-3,5)\) y \(\vec{\mathbf{d}}=(7,-4,1,11)\).
  Mostrar que \(\vec{\mathbf{ab}}\) es paralelo a \(\vec{\mathbf{cd}}\).
\item
  Sea \(\vec{\mathbf{a}}=(3,-1,4)\) y \(\vec{\mathbf{b}}=(4,2,-1)\).
  Encontrar
  \(a)\quad\|\vec{\mathbf{a}} + \vec{\mathbf{b}}\|\quad\quad\quad b)\quad \|\vec{\mathbf{a}} - \vec{\mathbf{b}}\|\quad\quad\quad c)\quad \|2\vec{\mathbf{a}}-\vec{\mathbf{b}}\|\quad\quad\quad d)\quad \|\frac{1}{2}\vec{\mathbf{a}}+\frac{1}{4}\vec{\mathbf{b}}\|\).
\item
  Encontrar un vector \(\vec{\mathbf{x}}\),
  \(\vec{\mathbf{x}}\ne \vec{\mathbf{0}}\) que sea ortogonal a
  \(\vec{v}=(1,3,2)\). Encuentra una condición que cumplan todos los
  vectores ortogonales al vector \(\vec{v}\).
\item
  Sean \(\vec{A}=(5,3,3)\), \(B=(1,3,1)\) y \(C=(2,6,-1)\). Uno de los
  ángulos del tríangulo \(ABC\) es recto. ¿qué ángulo es?
\end{enumerate}

\hypertarget{producto-escalar-en-mathbbrn}{%
\section{\texorpdfstring{Producto escalar en
\(\mathbb{R}^n\)}{Producto escalar en \textbackslash mathbb\{R\}\^{}n}}\label{producto-escalar-en-mathbbrn}}

\begin{enumerate}
\def\labelenumi{\arabic{enumi}.}
\item
  Probar usando la definición de la norma euclidiana que para todo
  \(c\in\mathbb{R}\) y \(\vec{x}\in \mathbb{R}^n\),
  \(\|c\cdot \vec{x}\|=|c|\|\vec{x}\|\).
\item
  Probar que
  \((a\vec{v}+b\vec{w})\boldsymbol{\cdot}(c\vec{v}+d\vec{w}) = ac\|\vec{v}\|^2 + (ad+bc)(\vec{v}\boldsymbol{\cdot}\vec{w})+bd\|\vec{w}\|^2\),
  para todo \(a,b,c,d\in \mathbb{R}\) y
  \(\vec{v},\vec{w}\in\mathbb{R}^n\)
\item
  Probar la identidad de \emph{polarización}. Para todo
  \(\vec{v},\vec{w}\in\mathbb{R}^n\) \[
  \vec{v}\boldsymbol{\cdot}\vec{w}=\frac{1}{4}(\|\vec{v}+\vec{w}\|^2-\|\vec{v}-\vec{w}\|^2)
  \]
\item
  Encontrar \(a\) tal que \(\vec{v}=(2,a,-3)\) sea ortogonal a
  \(\vec{w}=(-1,3,-2)\)
\item
  Escribir una condición para caracterizar a todos los vectores
  \(\vec{\mathbf{x}}\in\mathbb{R}^{3}\) que sean ortogonales a los
  vectores \(\vec{\mathbf{a}}=(1,3,2)\), \(\vec{\mathbf{b}}=(-1,-2,1)\)
  y \(\vec{\mathbf{c}}=(0,1,4)\).
\item
  Si \(A=(1,-2,2)\) y \(B=(7,-5,4)\) si \(\vec{M}\) es el punto medio de
  \(\vec{OA}\) y \(\vec{N}\) es el punto medio de \(\vec{OB}\), mostrar
  que \(d(\vec{M},\vec{N})=\frac{1}{2}d(\vec{A},\vec{C})\)
\item
  Usando las propiedades del producto punto y de la norma euclidiana
  simplificar
\end{enumerate}

\[
a)\,\,  \|\big(\vec{a}\bullet\vec{b}\big) \vec{c} \|
\quad\quad
b)\,\,  \|\big(\big(\vec{a}\bullet\vec{b}\big)a\big) \times \vec{c} \|
\quad\quad
c) \,\, \|\vec{a}\times \big(\vec{b} \times \vec{c}\big)\|
\]

\begin{enumerate}
\def\labelenumi{\arabic{enumi}.}
\tightlist
\item
  Encuentre el área (usando las proyecciones ortogonales) y perímetro
  del triángulo con vértices
\end{enumerate}

\begin{enumerate}
\def\labelenumi{\alph{enumi})}
\item
  \((0,0),(3,1),(7,0)\)
\item
  \((1,3),(2,1),(7,2)\)
\end{enumerate}

\begin{enumerate}
\def\labelenumi{\arabic{enumi}.}
\tightlist
\item
  Sea \(\vec{r}(t)=(\cos\omega t, \sin \omega t)\).
\end{enumerate}

\begin{itemize}
\item
  Calcular \(\frac{d\vec{r}}{dt}(t)\).
\item
  Demostrar que el vector velocidad y el vector desplazamiento
  \(\vec{r}(t)\) son ortogonales.
\end{itemize}

\begin{enumerate}
\def\labelenumi{\arabic{enumi}.}
\item
  Calcular los ángulos que el vector \(\vec{v}=(1,-1,0,\sqrt{2})\) con
  \(\vec{e}_1=(1,0,0,0)\), \(\vec{e}_2=(0,1,0,0)\),
  \(\vec{e}_3=(0,0,1,0)\) y \(\vec{e}_4=(0,0,0,1)\).
\item
  Demostrar las propiedades del producto cruz
\end{enumerate}

\begin{itemize}
\item
  Para todo \(\vec{u},\vec{v}\in\mathbb{R}^3\) se cumple que
  \(\vec{u} \times \vec{v}=- \vec{v} \times \vec{u}\)
\item
  Demostrar que \(\vec{u}\times \vec{v}\) es un vector ortogonal a
  \(\vec{u}\) y a \(\vec{v}\).
\item
  Demostrar que si \(\vec{u},\vec{v},\vec{w}\in\mathbb{R}^3\) se cumple
  que
  \(\vec{u} \times (\vec{v}+\vec{w})= \vec{u}\times \vec{v} + \vec{u}\times \vec{w}\).
\end{itemize}

\hypertarget{rectas-y-planos}{%
\section{Rectas y planos}\label{rectas-y-planos}}

\begin{enumerate}
\def\labelenumi{\arabic{enumi}.}
\item
  Encontrar las ecuaciones vectorial y paramétricas del plano que pasa
  por \(A=(-1,2,2)\), \(B=(0,3,1)\) y \(C=(0,0,2)\).
\item
  Determinar la ecuación del plano que pasa por \(A=(1,0,4)\) que es
  ortogonal a la línea que pasa por el origen y \((2,3,4)\)
\item
  Escribir la ecuación del plano que pasa por el punto \((3,7,2)\) y es
  paralelo a las rectas \(l_1:\)
  \(\frac{x-1}{4}=\frac{y+2}{1}=\frac{z-4}{2}\) y \(l_2:\)
  \(\frac{x-3}{5}=\frac{y-2}{3}=\frac{z+1}{1}\)
\item
  Escribir la ec. del plano que pasa por \((1,0,-1)\) y contiene a la
  recta \(\frac{x}{2}=\frac{y+1}{3}=\frac{z-2}{1}\).
\item
  Sea \(E\) el conjunto de puntos de la forma \[
    \{ \vec{\mathbf{v}} \in \mathbb{R}^{3} \,:\, \vec{\mathbf{v}} = (3,2,1) + s (2,0,1) + t(1,1,2), s,t\in\mathbb{R}   \}
    \] Mostrar que \(E\) es un plano y encontrar su ecuación cartesiana.
\item
  Sea \(E\) el plano con ecuación \(x+y-2z=3\). Si \(l\) es la recta con
  ecuación \(X(t)=(2,1,3) + t(0,1,4)\). Hallar la intersección de la
  recta y el plano.
\item
  Sea \(\vec{A}=(-1,2,3)\), \(\vec{B}=(2,5,-3)\), \(\vec{C}=(4,1,-1)\) y
  \(\vec{D}=(1,3,-3)\).
\end{enumerate}

\begin{enumerate}
\def\labelenumi{\alph{enumi})}
\item
  Encuentre el vector de los puntos de trisección de la recta que pasa
  por \(\vec{A}\) y \(\vec{B}\).
\item
  Encuentre la distancia entre los puntos \(A\) y \(B\).
\item
  Encontrar la ecuación del plano que pasa por \(\vec{A}\) y
  perpendicular a \(\vec{OB}\).
\item
  Encontrar la ecuación del plano que pasa por el origen \(\textbf{O}\)
  y los puntos \(\vec{A}\) y \(\vec{B}\).
\item
  Encontrar la distancia de \(\vec{OD}\) al plano formado por
  \(A,B\,\mbox{ y } C\). \emph{Sugerencia}: revisar la fórmula o bien
  calcular la recta que pasa por \(OD\) y dirección \(\vec{d}=\vec{n}\).
\item
  Encontrar la ecuación vectorial (paramétrica) que pasa por \(\vec{A}\)
  y paralelo a \(\vec{BC}\).
\item
  Encontrar la ecuación de la recta que pasa por \(\vec{D}\) y forma un
  ánuglo de \(30°\) con la normal al plano formado por \(ABC\).
\end{enumerate}

\begin{enumerate}
\def\labelenumi{\arabic{enumi}.}
\setcounter{enumi}{8}
\tightlist
\item
  Se dan los vértices de un triángulo en \(A=(0,0,0)\), \(B=(0,4,4)\) y
  \(C=(3,3,0)\).
\end{enumerate}

\begin{enumerate}
\def\labelenumi{\roman{enumi})}
\item
  Calcular los ángulos internos en cada vértice y el perímetro del
  triángulo
\item
  Calcular el plano que contiene al triángulo.
\item
  Encontrar el conjunto intersección del plano en ii) y el plano
  \(2x-y-2z=4\)
\end{enumerate}

\begin{enumerate}
\def\labelenumi{\arabic{enumi}.}
\setcounter{enumi}{9}
\item
  Para determinar la distancia \(d\), mas corta entre 2 rectas \(l_1\) y
  \(l_2\) si \(\vec{d}_1\) y \(\vec{d}_2\) son direcciones de \(l_1\) y
  \(l_2\) respectivamente y \(\vec{n}=\vec{d_1} \times \vec{d_2}\).
  Suponga que \(\vec{PR}\) une a un punto de la recta \(l_1\) y a un
  punto de la recta \(l_2\). Usando la fórmula de la longitud de la
  proyección ortogonal, demuestre que la distancia \(d\) entre las
  rectas está dada por \[
  d = \frac{|\vec{PR} \cdot \vec{n}|}{\|\vec{n}\|}
  \]
\item
  Use el ejercicio anterior para calcular lo siguiente
\end{enumerate}

\begin{enumerate}
\def\labelenumi{\alph{enumi})}
\item
  Determinar la distancia más corta entre la recta \(l_1\) que pasa por
  \(P(1,2,-1)\) y \(Q(0,-1,2)\) y la recta que pasa por \(R(-1,2,0)\) y
  \(S(-1,0,-2)\).
\item
  Determinar la distancia más corta entre las 2 rectas cuyas ecuaciones
  paramétricas son \[
  \begin{array}{cccc}
  x =& 5-2t & \quad\quad &x = 5-s \\
  y =&-4+2t & \quad\quad &y = -2-2s\\
  z =& 2-t  & \quad\quad &z = 1+s 
  \end{array}
  \]
\item
  Encuentre la distancia más corta entre las rectas \[
  \begin{array}{cc}
  l_1: & \frac{x-2}{3} = \frac{y-5}{2} = \frac{z-1}{-1} \\
  l_2: & \frac{x-4}{-4} = \frac{y-5}{4} = \frac{z+2}{1} 
  \end{array}
  \]
\end{enumerate}

\begin{enumerate}
\def\labelenumi{\arabic{enumi}.}
\setcounter{enumi}{11}
\tightlist
\item
  Encuentre la distancia del punto dado al plano
\end{enumerate}

\begin{enumerate}
\def\labelenumi{\alph{enumi})}
\item
  \(P(4,0,1)\) y \(\Pi_1: 2x-y+8z=3\)
\item
  \(P(-7,-2,-1)\) y \(\Pi_2: -2x+8z=-5\)
\end{enumerate}

\newpage

\hypertarget{sistemas-de-ecuaciones}{%
\section{Sistemas de ecuaciones}\label{sistemas-de-ecuaciones}}

\begin{enumerate}
\def\labelenumi{\arabic{enumi}.}
\item
  Considere los siguientes sistemas de ecuaciones lineales: \[
  \begin{array}{cc}
    2x_1 + 3x_2 &= 5 \\
    x_1  - 4x_2 &= -3 \\
    2x_1 - 8x_2 & = -6
  \end{array} 
  \quad \quad \quad
  \begin{array}{cc}
    2x_1 + 3x_2 &= 5 \\
    x_1  - 4x_2 &= -3 
  \end{array}
  \] Determine el conjunto solución de cada uno de los sistemas ¿Son
  sistemas equivalentes?
\item
  Encontrar todos los vectores en \(\mathbb{R}^3\) que sean ortogonales
  a \((1,2,3)\) y a \((-2,0,1)\).
\item
  Usar eliminación gaussiana para resolver los siguientes sistemas
\end{enumerate}

\begin{multicols}{2}
\begin{equation*}
  \begin{array}{cccc}
     x_1   & + x_2  & + x_3  & = 1  \\
     x_1   & + 2x_2 & +2x_3  & = 1  \\
     x_1   & + 2x_2 & +3x_3  & = 1
  \end{array}
\end{equation*}

\begin{equation*}
  \begin{array}{cccc}
     2x_1  & - x_2  &       & = 0  \\
     -x_1  & + 2x_2 & -x_3  & = 0  \\
           & - x_2 & + x_3  & = 1
  \end{array}
\end{equation*}
\begin{equation*}
  \begin{array}{cccc}
            &   4x_2  & -3 x_3  & = 3  \\
     -x_1   & + 7x_2  & -5x_3   & = 4  \\
     -x_1   & + 8x_2  & -6x_3   & = 5
  \end{array}
\end{equation*}

\begin{equation*}
  \begin{array}{ccccc}
     x_1   & + x_2  & +x_3  &+x_4   & = 1  \\
     x_1   & + x_2  & +3x_3 &+3x_4  & = 3  \\
     x_1   & + x_2  & +2x_3 &+3x_4  & = 3  \\
     x_1   & +3x_2  & +3x_3 &+3x_4  & = 4
  \end{array}
\end{equation*}
\end{multicols}

\begin{enumerate}
\def\labelenumi{\arabic{enumi}.}
\setcounter{enumi}{3}
\tightlist
\item
  Determinar el valor de \(\lambda\) para que el siguiente sistema de
  ecuaciones tenga soluciones no nulas (y de hecho una infinidad).
\end{enumerate}

\begin{equation*}
  \begin{array}{cccc}
     x_1   & + 2x_2   & -2 x_3        & = 0  \\
     2x_1  & - x_2    & +\lambda x_3  & = 0\\
     3x_1  & + x_2    & -x_3z           & = 0
  \end{array}
\end{equation*}

\begin{equation*}
  \begin{array}{cccc}
     x_1   &          & + x_3  & = 4  \\
     2x_1  & + x_2    & +3x_3  & = 5\\
    -3x_1  & -3x_2    & +(\lambda^2-5\lambda)x_3         & = \lambda-8
  \end{array}
\end{equation*}

\begin{enumerate}
\def\labelenumi{\arabic{enumi}.}
\setcounter{enumi}{4}
\tightlist
\item
  Encontrar las soluciones del sistema de ecuaciones aplicando
  Gauss-Jordan
\end{enumerate}

\begin{multicols}{2}
\begin{equation}
  \begin{array}{cccc}
     x_1   & + x_2   & +x_5    &= 0  \\
     x_1   & + x_2   & -x_3    & = 0\\
     x_3   & + x_4   & +x_5    & = 0
  \end{array}
\end{equation}\quad

\begin{equation}
  \begin{array}{cccc}
     x_1   & + x_2   & -x_3    &= 2  \\
     x_1   & -3 x_2   & +3x_3    & = 0\\
     2x_1   & + x_2   & +x_3    & = 6
  \end{array}
\end{equation}
\begin{equation}
  \begin{array}{cccc}
     x_1   & + x_2   & -x_3    &= 0  \\
     2x_1  & +2 x_2   & -x_3    & = 2\\
     -x_1   & - x_2   & +x_3    & = 0
  \end{array}
\end{equation}

\begin{equation}
  \begin{array}{ccccl}
     x_1   & + x_2    & -2x_3 & +x_4    &= 3  \\
     2x_1  & +2 x_2   & +x_3  &-3x_4    & = 1\\
     x_1   & - 2x_2    & +x_3 &+x_4     & = 3 \\
     x_1   & - x_2    & -x_3 &+2x_4     & = 4 
  \end{array}
\end{equation}
\end{multicols}

6.Determinar el valor de \(\lambda\) para que el siguiente sistema de
ecuaciones tenga una infinidad de soluciones. ¿En que casos no tiene
solución?

\begin{equation*}
  \begin{array}{cccccc}
     x_1        & + & 2x_2  & - & 2 x_3            & = 0  \\
    2x_1        & - &  x_2  & + &  \lambda x_3     & = 0\\
    3x_1        & + &  x_2  & - & \lambda  x_3     & = 0
  \end{array}
\end{equation*}

7.Para las siguientes matrices:

\begin{itemize}
\item Usando eliminación gaussiana encontrar su forma escalonada reducida. 
\item Escribir el sistema de ecuaciones resultante $Ax=0$.
\item Resolver el sistema de ecuaciones y expresar el conjunto solución.
\end{itemize}

\begin{equation*}
\begin{array}{ccccc}
A =& \begin{bmatrix} 
      2  & 3 & 6 \\ 
      4  & 7 & 9 \\
      3  & 5 & 4  
     \end{bmatrix} & \,\,\, &
B= & \begin{bmatrix} 
      1   & -2 & 1 & 0\\ 
      -2  &  4 & 0 & 2\\
      -1  &  2 & 1 & 2  
     \end{bmatrix}
\end{array}
\end{equation*}

\begin{equation*}
\begin{array}{ccccc}
C = & \begin{bmatrix} 
      1  &-2 & 1 & 0 \\ 
     -2  & 4 & 0 & 2\\
     -1  & 2 & 1 & 2  
     \end{bmatrix} & \,\,\, &
D = & \begin{bmatrix} 
      1   &  0 & 3 \\ 
      0   &  1 & -1\\
      -1  &  1 & -4  
     \end{bmatrix}
\end{array}
\end{equation*}

\begin{enumerate}
\def\labelenumi{\arabic{enumi}.}
\setcounter{enumi}{7}
\tightlist
\item
  Determine todas las soluciones del sistema de ecuaciones aplicando el
  método de Gauss-Jordan
\end{enumerate}

\begin{multicols}{2}
\begin{equation}
  \begin{array}{cccc}
     x_1   & + x_2   & +x_3    & -2x_4    &= 3  \\
    2x_1   & + x_2   & +3x_3   & +2x_4    &= 5\\
           & - x_2   & +x_3    & +6x_4    &= 3
  \end{array}
\end{equation}\quad

\begin{equation}
  \begin{array}{cccc}
     x_1   & + x_2   & -x_3    &= 2  \\
     x_1   & -3 x_2   & +3x_3    & = 0\\
     2x_1   & + x_2   & +x_3    & = 6
  \end{array}
\end{equation}
\begin{equation}
  \begin{array}{cccc}
     x   & - y   & +z    & = 2  \\
     3x  & + y   & -z    & = 3\\
     x   & + y   & +z    & = 3
  \end{array}
\end{equation} \quad

\begin{equation}
  \begin{array}{cccc}
     x   & - y   & +3z    & = 2  \\
     2x  & - y   & +z    & = 3\\
     x   & + y   & +z    & = 4
  \end{array}
\end{equation}

\begin{equation}
  \begin{array}{cccc}
     x   & +2 y   & +z    & = 24  \\
     2x  & -3 y   & +z    & = -1  \\
     x   & -2 y   & +2z   & = 7
  \end{array}
\end{equation}

\begin{equation}
  \begin{array}{cccc}
     a   & -2 b   & -4c    & = -3  \\
     2x  & -3 y   & +z    & = -1  \\
     x   & -2 y   & +2z   & = 7
  \end{array}
\end{equation}
\end{multicols}

\hypertarget{matrices}{%
\section{Matrices}\label{matrices}}

\begin{enumerate}
\def\labelenumi{\arabic{enumi}.}
\item
  Calcula los productos \(ABC\) y \(BA\) de las matrices, con \[
  \begin{array}{lcr}
  A=\begin{pmatrix}
   991 & 992 & 993 \\
   994 & 995 & 996 \\
   997 & 998 & 999
    \end{pmatrix}, \quad &
  B = \begin{pmatrix}
   12 & -6 & -2 \\
   18 & -9 & -3 \\
   24 &-12 & -4
   \end{pmatrix}, \quad &
  C = \begin{pmatrix}
   1 & 1  \\
   1 & 2  \\
   3 & 0 
   \end{pmatrix}
  \end{array}
  \]
\item
  \begin{enumerate}
  \def\labelenumii{\alph{enumii})}
  \tightlist
  \item
    Una matriz se dice \emph{idempotente} si \(A^2 = A\). Probar que \[
    B= \begin{pmatrix}
     2 & -3 & -5 \\
    -1 &  4 & 5  \\
     1 & -3 &- 4
    \end{pmatrix}
    \] es idempotente
  \end{enumerate}
\end{enumerate}

\begin{enumerate}
\def\labelenumi{\alph{enumi})}
\setcounter{enumi}{1}
\tightlist
\item
  Demuestre qu si \(A\) es idempotente, entonces \(B=I_n-A\) también es
  idempotente,
\item
  Si \(A\) y \(B\) son como en b) entonces demuestre que
  \(AB= BA = \mathbf{0}\)
\end{enumerate}

\begin{enumerate}
\def\labelenumi{\arabic{enumi}.}
\setcounter{enumi}{1}
\tightlist
\item
  Se dice que una matriz \(n\times n\) es \emph{involutiva} si y sólo si
  \(A^2 = I_n\).
\end{enumerate}

\begin{enumerate}
\def\labelenumi{\alph{enumi})}
\item
  Verifica que
  \(A=\begin{pmatrix} 0 & 1 & -1 \\ 4 & -3 & 4 \\ 3 & -3 & 4 \end{pmatrix}\)
  es involutiva
\item
  Demuestre que si \(A\) es involutiva entonces \(\frac{1}{2}(I_n+A)\) y
  \(\frac{1}{2}(I_n-A)\) son idempotentes y su producto es igual a la
  matriz cero.
\end{enumerate}

\begin{enumerate}
\def\labelenumi{\arabic{enumi}.}
\setcounter{enumi}{2}
\tightlist
\item
  Para las siguientes matrices, encontrar la fórmula para \(A^{n}\) con
  \(n\in\mathbb{N}\),
\end{enumerate}

\[
\begin{array}{cr}
 A= \begin{pmatrix}
    \cos x       & -\mbox{sen} x \\
    \mbox{sen} x & \cos x
    \end{pmatrix}, &
A= \begin{pmatrix}
    a       & 1 \\
    0       & a
    \end{pmatrix} 
\end{array}
\]

\begin{enumerate}
\def\labelenumi{\arabic{enumi}.}
\setcounter{enumi}{3}
\tightlist
\item
  En \(\mathbb{R}^3\) la matriz \[
   R_z(\phi)= \begin{pmatrix}
           \cos \phi       & -\mbox{sen} \phi & 0 \\
           \mbox{sen}\phi  & \cos \phi        & 0 \\
           0               &      0           & 1
   \end{pmatrix},
  \]
\end{enumerate}

\begin{enumerate}
\def\labelenumi{\alph{enumi})}
\item
  Calcular \(R_z(\phi)\vec{x}\) y exlicar rota al vector \(\vec{x}\)
  alrededor del eje \(z\).)
\item
  Demuestre que si \(\vec{v}' = R_z(\phi)\vec{v}\) entonces
  \(\|\vec{v´}\|=\|\vec{v}\|\)
\end{enumerate}

\begin{enumerate}
\def\labelenumi{\arabic{enumi}.}
\setcounter{enumi}{4}
\item
  Dada la matriz
  \(\mathbf{X} =\begin{pmatrix} 0 & 2 \\ -1 & 0 \end{pmatrix}\)
  encontrar todas las matrices \(\mathbf{Y}\) tales que
  \(\mathbf{X}\mathbf{Y} = \mathbf{Y}\mathbf{X}\)
\item
  El \emph{conmutador} de dos matrices cuadradas \(A\) y \(B\) se define
  como \[
  \big[A,B\big] = AB - BA
  \]
\end{enumerate}

\begin{enumerate}
\def\labelenumi{\alph{enumi})}
\item
  Mostrar que las matrices conmutan con la multiplicación si
  \(\big[A,B\big] = \mathbf{0}\)
\item
  Calcular el conmutador de las siguientes matrices
\end{enumerate}

\begin{itemize}
\item
  \(A=\begin{pmatrix} 1 & 0 \\1 & -1\end{pmatrix},\quad B=\begin{pmatrix} 2 & 1 \\-2 & 0\end{pmatrix}\)
\item
  \(A=\begin{pmatrix} -1 & 3 \\ 3 & -1\end{pmatrix},\quad B=\begin{pmatrix} 1 & 7 \\-2 & 0\end{pmatrix}\)
\end{itemize}

\begin{enumerate}
\def\labelenumi{\alph{enumi})}
\setcounter{enumi}{2}
\tightlist
\item
  Demostrar que \(\big[A,I\big] = \mathbf{0}\) para toda \(A\) matriz
  \(n\times n\)
\end{enumerate}

\begin{enumerate}
\def\labelenumi{\arabic{enumi}.}
\setcounter{enumi}{6}
\item
  Calcular \(M^2\), \(M^3\) si \(M\) está dada por \[
  M= \begin{pmatrix}
            0   & 0 & 0 & 0 \\
           3    & 0 & 0 & 0 \\
           1    & -1& 0 & 0 \\
           -1   & 2& 1 & 0   
   \end{pmatrix}, \quad
  M= \begin{pmatrix}
            1   & 0 & 0 & 0 \\
           3    & 1 & 0 & 0 \\
           1    & -1& 1 & 0 \\
           -1   & 2& 1 & 1   
   \end{pmatrix}
  \]
\item
  Si \(A=\begin{pmatrix} 1 & 2 \\ -1 & 0 \end{pmatrix}\) calcular
  \(A^5\).
\item
  Calcular \((ABC)^{-1}\) aplicando las propiedades de las matrices
  inversas \[
  A= \begin{pmatrix} 1 & 2 \\ 0 & 1 \end{pmatrix},\quad\quad B= \begin{pmatrix} 2 & 5 \\ 3 & 7 \end{pmatrix}\quad \mbox{ y } \quad 
  C= \begin{pmatrix} 3 & 2 \\ 4 & 2\end{pmatrix}
  \]
\item
  Si \(A=\begin{pmatrix} 0 & 1 \\ -1 & 0 \end{pmatrix}\) calcular
  \(A^{1000}\).
\item
  Si \(\begin{pmatrix} 1 & 2 \\ 2 & -1 \end{pmatrix}\) calcular
  \(B^{1000}\).
\item
  Calcular \(A=\begin{pmatrix} 1 & 1 \\ 0 & 1 \end{pmatrix}\)
\end{enumerate}

\begin{itemize}
\item
  calcular \(A^2\).
\item
  calcular \(A^3\).
\item
  calcular \(A^{k}\). \emph{observe} \emph{un} \emph{patrón} \emph{y}
  \emph{aplique} \emph{inducción} \emph{matemática}.
\end{itemize}

\begin{enumerate}
\def\labelenumi{\arabic{enumi}.}
\item
  Para \(\begin{pmatrix} 1/2 & \alpha \\ 0 & 1/2 \end{pmatrix}\)
  determinar para que casos de \(\alpha\) existe el límite
  \(\lim_{n\rightarrow \infty} A^{n}\). (\emph{sugerencia}:
  \emph{determinar} \emph{algunas} \emph{potencias} \emph{de} \(A^n\))
\item
  Sea \(\vec{\mathbf{e}}_j\) el vector unitario \(j\), el cual contiene
  en la j-ésima posición a 1 y en las demás coordenadas es igual a 0.
  Para una matrix \(A\) describir los siguientes productos
\end{enumerate}

\begin{enumerate}
\def\labelenumi{\alph{enumi})}
\tightlist
\item
  \(A\cdot e_j\) b) \(e_i^T\cdot A\) c) \(e_i^T\cdot A\cdot e_j\)
\end{enumerate}

\hypertarget{matrices-inversas}{%
\section{Matrices inversas}\label{matrices-inversas}}

\begin{enumerate}
\def\labelenumi{\arabic{enumi}.}
\tightlist
\item
  Encuentre la inversa de la matriz
\end{enumerate}

\begin{equation}
  \begin{pmatrix}
    1 & 2 & 3 & 4 \\
    0 & 2 & 3 & 4 \\
    0 & 0 & 3 & 4 \\
    0 & 0 & 0 & 4
  \end{pmatrix}
\end{equation}

\begin{enumerate}
\def\labelenumi{\arabic{enumi}.}
\setcounter{enumi}{1}
\tightlist
\item
  Generalize el problema para
\end{enumerate}

\begin{equation}
  A_5 = \begin{pmatrix}
          1 & 2 & 3 & 4 & 5 \\
          0 & 2 & 3 & 4 & 5 \\
          0 & 0 & 3 & 4 & 5 \\
          0 & 0 & 0 & 4 & 5 \\
          0 & 0 & 0 & 0 & 5 \\
        \end{pmatrix}
\end{equation}

\begin{enumerate}
\def\labelenumi{\arabic{enumi}.}
\setcounter{enumi}{2}
\tightlist
\item
  Encontrar la inversa de matriz
\end{enumerate}

\begin{equation}
  A_5 = \begin{pmatrix}
          1           & \frac{1}{2} & \frac{1}{3} & \frac{1}{4} \\
          \frac{1}{2} & \frac{1}{3} & \frac{1}{4} & \frac{1}{5} \\
          \frac{1}{3} & \frac{1}{4} & \frac{1}{5} & \frac{1}{6} \\
          \frac{1}{4} & \frac{1}{5} & \frac{1}{6} & \frac{1}{7}
        \end{pmatrix}
\end{equation}

\begin{enumerate}
\def\labelenumi{\arabic{enumi}.}
\setcounter{enumi}{3}
\item
  Usando el método de Gauss-Jordan encontrar la inversa de
  \begin{equation}
    A= \begin{bmatrix} 
     1    & 0 & 0 & \ldots & 0 \\
     a_2  & 1 & 0 & \ldots & 0 \\
     a_3  & 0 & 1 & \ldots & 0 \\
   \vdots & \vdots & \ddots & \vdots \\
     a_{n-1} & 0 & 0 \ldots & 1
    \end{bmatrix}
  \end{equation}
\item
  Calcular la inversa \(A^{-1}\) de las siguientes matrices
\end{enumerate}

\begin{multicols}{2}
a) 
\begin{equation*}
\begin{array}{cc}
A= & 
  \begin{bmatrix} 
    -4 & -3\\ 
    -3 & 2
  \end{bmatrix}
\end{array}
\end{equation*}

b) 
\begin{equation*}
\begin{array}{cc}
A= & 
  \begin{bmatrix} 
    0  & 1 \\ 
    1  & 1
  \end{bmatrix}
\end{array}
\end{equation*}

c) 
\begin{equation*}
\begin{array}{cc}
A= & 
  \begin{bmatrix} 
    3  & 0 \\ 
    9  & 3
  \end{bmatrix}
\end{array}
\end{equation*}

d)
\begin{equation*}
\begin{array}{cc}
A= & 
  \begin{bmatrix} 
    2  & 1 & 2 \\ 
    0  & 3 & -1 \\
    4  & 1 & 1
  \end{bmatrix}
\end{array}
\end{equation*}

e) 

\begin{equation*}
\begin{array}{cc}
A= & 
  \begin{bmatrix} 
    3  & 1 & -1 \\ 
    2  & 1 &  0 \\
    1  & 2 & 4
  \end{bmatrix}
\end{array}
\end{equation*}
\end{multicols}

\begin{enumerate}
\def\labelenumi{\arabic{enumi}.}
\setcounter{enumi}{5}
\tightlist
\item
  Dadas las matrices \(A=\begin{bmatrix} 5 & 3 \\ 3 & 2 \end{bmatrix}\),
  \(B=\begin{bmatrix} 6 & 2 \\ 2 & 4 \end{bmatrix}\) y
  \(C=\begin{bmatrix} 4 & -2 \\ -6 & 3 \end{bmatrix}\) resuelva cada una
  de las siguientes ecuaciones de matrices:
\end{enumerate}

\begin{enumerate}
\def\labelenumi{\alph{enumi})}
\tightlist
\item
  \(AX + B = C\).
\item
  \(XA + B = C\).
\item
  \(XA + C = X\).
\end{enumerate}

\begin{enumerate}
\def\labelenumi{\arabic{enumi}.}
\setcounter{enumi}{6}
\tightlist
\item
  Calcular \(A+B\), \(AB\), \(A-B\), \(A^{-1}\), \(B^{-1}\) y
  \((AB)^{-1}\) si
\end{enumerate}

\begin{equation*}
\begin{array}{ccccc}
A= & 
  \begin{bmatrix} 
    1  & 0 \\ 
    1  & 1
  \end{bmatrix}   & \,\,\,\,\,\, & 
B = &  \begin{bmatrix} 
    1  & 1 \\ 
    0  & 1
  \end{bmatrix} 
\end{array}
\end{equation*}

\begin{enumerate}
\def\labelenumi{\arabic{enumi}.}
\setcounter{enumi}{7}
\tightlist
\item
  Sean \begin{equation*}
  \begin{array}{cccccccc}
  A= & 
    \begin{bmatrix} 
   2  & 1 & 2 \\ 
   0  & 3 & -1 \\
   4  & 1 & 1
    \end{bmatrix}   & \,\,\,\,\,\, & 
  B = &  \begin{bmatrix} 
   3   \\ 
   1   \\
   -1  
    \end{bmatrix} & \,\,\,\,\,\, &
  C= & \begin{bmatrix} 
   1  & -1 & 1
    \end{bmatrix}
  \end{array}
  \end{equation*}
\end{enumerate}

Calcular \(AB\), \(CAB\), \(B^{T}A^{T}C^{T}\), \(CB\), \(BC\),
\(AC^{T}B\), \(CBA\).

\begin{enumerate}
\def\labelenumi{\arabic{enumi}.}
\setcounter{enumi}{8}
\item
  Considere la siguiente matriz \(A\) ¿para qué valores de \(k\), la
  matriz es invertible? \begin{equation*}
    \begin{bmatrix}
    k+3 & -1    &  1  \\
    5   & k-3   &  1\\
    6   & -6    & k+4
    \end{bmatrix}
  \end{equation*}
\item
  Aplicando el método de Gauss-Jordan, mostrar que las siguientes
  matrices son no invertibles \begin{equation*}
  \begin{array}{cccc}
  A= & 
    \begin{bmatrix} 
  4 & -3\\ 
  -8 & 6
    \end{bmatrix} \quad \quad
     &
  A= & 
    \begin{bmatrix} 
  0  & 1 \\ 
  1  & 1
    \end{bmatrix}
  \end{array}
  \end{equation*}
\item
  Mostrar que para todo real \(\theta\) la matriz es invertible y
  encontrar su inversa. \begin{equation*}
  \begin{array}{cc}
  A= & 
    \begin{bmatrix} 
  \sin(\theta)  & \cos(\theta)  & 0\\ 
  \cos(\theta)  & -\sin(\theta)  & 0\\
  0  & 0  & 1
    \end{bmatrix}
  \end{array}
  \end{equation*}
\item
  ¿Para que valores de \(\lambda\) el sistema tiene solución única,
  infinidad de soluciones o ninguna solución?
\end{enumerate}

\begin{multicols}{2}
\begin{equation*}
  \begin{array}{cccc}
     \lambda x   & + y   & + z    &= 1  \\
      x   & + \lambda y  & + z    & = 1\\
      x   & +  y  & +\lambda z    & =1
  \end{array}
\end{equation*}
\begin{equation*}
  \begin{array}{cccc}
     (1+\lambda) x   & + y   & + z    & = 1  \\
      x   & +(1+\lambda) y  & + z     & = \lambda  \\
      x   & +  y  & + (1+\lambda)z    & = \lambda^2
  \end{array}
\end{equation*}
\end{multicols}

\begin{enumerate}
\def\labelenumi{\arabic{enumi}.}
\setcounter{enumi}{12}
\item
  Use el método de Gauss-Jordan para encontrar la inversa de la matriz o
  explicar por que no existe. \begin{equation*}
  \begin{array}{ccc}
    \begin{bmatrix} 
  1  & -1 & 2 \\ 
  3  & 1 & 2 \\
  2  & 3 & -1
    \end{bmatrix} & \,\,\,\,\,\, &
    \begin{bmatrix} 
  2  & 3 & 0 \\ 
  1  &-2 &-1 \\
  2  & 0 &-1
    \end{bmatrix} \\
  \begin{bmatrix} 
  1  & 1 & 0 \\ 
  1  & 0 & 1 \\
  0  & 1 & 1
    \end{bmatrix} & \,\,\,\,\,\, &
    \begin{bmatrix} 
  a  & 0 & 0 \\ 
  1  & a & 0 \\
  0  & 1 & a
    \end{bmatrix}
  \end{array}
  \end{equation*}
\item
  Demostrar que si \(A\) es invertible y simérica entonces \(A^{-1}\)
  también es simétrica.
\item
  Si \(A\) es una matriz cuadrada de tamaño \(n\), tal que
  \(\mathbf{I}-A\) es invertible, probar que \[
  A(\mathbf{I}-A)^{-1} = (\mathbf{I}-A)^{-1}A
  \] \newpage
\end{enumerate}

\hypertarget{regla-de-cramer}{%
\section{Regla de Cramer}\label{regla-de-cramer}}

\begin{enumerate}
\def\labelenumi{\arabic{enumi}.}
\tightlist
\item
  Usa la regla de Cramer para resolver los siguientes sistemas de
  ecuaciones:
\end{enumerate}

\begin{multicols}{2}
\begin{equation*}
\begin{array}{cccccc}
x & + y &+ z &- w & = & 2  \\
  &   y &- z &+w  & = & 1  \\
  &     &  z &-w  & = & 0  \\
  &     &  z &+2w & = & 3
\end{array}
\end{equation*}

\begin{equation*}
\begin{array}{ccccc}
x & +y & +z & = & a \\
  & y  & +z & = & 0 \\
  &    & z  & = & 1 \\
\end{array}
\end{equation*}

\begin{equation*}
\begin{array}{ccccc}
x & + y &+ z  & = & 11  \\
2x& -6y &- z  & = & 0  \\
3x& +4y &+2z  & = & 0  
\end{array}
\end{equation*}

\begin{equation*}
\begin{array}{ccccc}
x & +2y & +3z & = & 7 \\
2x& +y  & +2z & = & 0 \\
-2x&+y  & -z  & = & 4 
\end{array}
\end{equation*}
\end{multicols}

\begin{enumerate}
\def\labelenumi{\arabic{enumi}.}
\setcounter{enumi}{1}
\tightlist
\item
  Usando la fórmula de la adjunta calcula la inversa de las siguientes
  matrices:
\end{enumerate}

\begin{multicols}{2}
\begin{equation*}
\begin{array}{ccccc}
1) A=& 
  \begin{pmatrix} 
    -2       & 3      & 2           \\ 
     6       & 0      & 3           \\
     4       & 1      & -1     
  \end{pmatrix}
 & \,\,\,\,\, &
2) B=&  \begin{pmatrix} 
     \cos\theta       & 0      & -\mbox{sen}\theta    \\ 
       0              & 1      & 0           \\
    \mbox{sen}\theta  & 0      & \cos\theta     
  \end{pmatrix}
\end{array}
\end{equation*}

\begin{equation*}
\begin{array}{ccccc}
3) A=& 
  \begin{vmatrix} 
     1       & 2      & -1           \\ 
     0       & 1      & 3           \\
     2       & 3      & -3     
  \end{vmatrix}
 & \,\,\,\,\, &
4) A=&  \begin{vmatrix} 
     2       & 4      & -1           \\ 
     1       & -1     & 3           \\
     3       & 4      & -3   
  \end{vmatrix}
\end{array}
\end{equation*}
\end{multicols}

\begin{enumerate}
\def\labelenumi{\arabic{enumi}.}
\setcounter{enumi}{2}
\item
  Determinar las condiciones sobre \(a\) y \(b\) tal que el sistema
  \begin{equation*}
    \begin{array}{cccc}
    x   & + y   & -a z    &= 1  \\
    2 x & +3 y  & 2z      & = b\\
    x   & -y    & +bz     & =a
    \end{array}
  \end{equation*}
\item
  Aplicando la regla de Cramer, indica todos los valores de \(a\) tales
  que la solución del siguiente sistema \emph{no} sea única.
  \begin{equation*}
  \begin{array}{cccccc}
  (8-a)x & +2y      & +3z & +aw  & = & 0 \\
    x & +(9-a)y  & +4z & + aw & = & 0 \\
    x & +2y      & +3z & + aw & = & 0 \\
    x & +2y      & +3z & + aw & = & 0
  \end{array}
  \end{equation*}
\end{enumerate}

Para el siguiente sistema de ecuaciones lineales \begin{equation*}
  \begin{array}{cccccc}
       x_1 & -  & 2x_2 & +  &  ax_3    & = 1  \\
     3 x_1 & +  & 2x_2 & +  &  x_3    & = 2\\
     2 x_1 &    &      &  + &  ax_3    & = 3
  \end{array}
\end{equation*}

\begin{enumerate}
\def\labelenumi{\alph{enumi})}
\item
  Usando la regla de Cramer, resolver el sistema para a = 2
\item
  Encontrar los valores para los cuales el sistema tiene una infinidad
  de soluciones.
\end{enumerate}

\end{document}
